\documentclass[]{article}

\usepackage{amsmath}
\usepackage{multirow}
\usepackage{natbib}
\usepackage{graphicx} 


\begin{document}

In this paper, we investigated the optimal strategy for detecting the signature of elastic scattering between dark matter and baryons using the 21 cm forest. Such an interaction suppresses the formation of small-scale structure, providing a powerful observational test of dark matter microphysics. Our goal was to identify the most sensitive observable, disentangle the dark matter signal from astrophysical effects, and determine the ideal redshift window for future observations.

To achieve this, we employed the \texttt{HAYASHI} semi-analytic framework to model the statistics of 21 cm absorbers from $z=7$ to $15$. We incorporated the effects of dark matter-baryon scattering through two channels: a cutoff in the halo mass function, $M_{cut}$, representing the suppression of structure formation, and a thermal cooling of the intergalactic medium. Using the differential optical depth distribution, $dN/d\tau$, as our primary observable, we performed a Fisher matrix forecast that included a realistic, redshift-dependent model for the number of background radio sources to project the constraining power of future radio observatories.

Our results demonstrate that the signature of dark matter-baryon scattering is overwhelmingly dominated by the kinematic suppression of low-mass minihalos. The thermal cooling channel was found to have a negligible impact on the 21 cm absorption statistics, as the signal is already saturated in the cold, high-redshift intergalactic medium. The suppression of minihalos imprints a distinct fingerprint on the shape of the optical depth distribution, preferentially removing low-$\tau$ absorbers. This shape-based signature provides a powerful discriminant to distinguish the dark matter interaction from astrophysical uncertainties that tend to rescale the overall amplitude of the signal.

We learned that there is a fundamental trade-off between intrinsic physical sensitivity and statistical constraining power. While the signature of structure suppression becomes more pronounced at higher redshifts, the decreasing number of bright background radio sources weakens the statistical power of the measurement. Our Fisher analysis synthesized these competing effects, identifying an optimal observational window at $z \sim 8-10$. This window represents a strategic balance, maximizing the distinguishability of the signal while retaining sufficient statistical power for a robust detection.

In conclusion, the 21 cm forest is a uniquely powerful probe of the fundamental nature of dark matter. By targeting the characteristic suppression of low optical depth systems within the optimal redshift window of $z \sim 8-10$, future radio observatories can place constraints on the velocity-independent dark matter-baryon scattering cross-section that are four to five orders of magnitude more stringent than current limits from the Cosmic Microwave Background. This work establishes a clear observational pathway for leveraging the 21 cm forest to explore a vast and otherwise inaccessible region of the dark matter parameter space.

\end{document}
                